\section{规则对比与选用}\label{sec:compare}

\subsection{三套规则的横向对照}

\begin{table}[htbp]
\centering\footnotesize
\caption{三套规则对照}\label{tab:compare}
\begin{tabular}{@{}>{\raggedright\arraybackslash}p{2.9cm}>{\raggedright\arraybackslash}p{3.7cm}>{\raggedright\arraybackslash}p{3.7cm}>{\raggedright\arraybackslash}p{3.5cm}@{}}
\toprule
\textbf{维度} & \textbf{规则一 期货趋势} & \textbf{规则二 A 股反转} & \textbf{规则三 期货统计套利} \\
\midrule
信号口径 & 时序（单品种对自身历史） & 截面（全市场排序） & 截面（产业链内配对） \\
方向 & 趋势（$+s$） & 反转（$-s$） & 反转（价差偏离） \\
主窗口 & $L=250$ / 60 / 10 & $L=60$（备选 120） & $L=60$（备选 120） \\
调仓频率 & 5 交易日 & 21 交易日 & 事件驱动（阈值触发） \\
年化换手 & $8.8\sim35$ & $28.1$ & $\sim27$（论文参照值） \\
成本占毛利 & $10.2\%\sim29.4\%$ & $31.9\%$ & $\sim6.3\%$（论文参照值） \\
波动率目标化 & \textbf{用} & 不用 & 不用 \\
入场门槛 & 信号强度 $>1\sigma$ & 无（排序） & \textbf{协整检验 + 半衰期筛选} \\
因子中性化 & 不做 & \textbf{行业 + 市值} & 配对天然对冲 \\
主要失效模式 & 慢性出血 + 崩溃跳崖 & 不收敛 + 资金断裂 & 不收敛 + 保证金挤兑 \\
主要防线 & 恐慌状态 + 拥挤监控 & 单票止损 + 现金冗余 & 止损 + 保证金冗余 \\
容量约束来源 & 拥挤（信号同质化） & 成本 & 成本 + 产业链深度 \\
论文支持的置信度 & 中（表 5 时序口径直接支持） & 高（表 6 直接支持） & 中（方向支持，实现改造未经检验） \\
\bottomrule
\end{tabular}
\end{table}

\emph{注}：表中“主窗口”一行的取值沿用论文口径；第~\ref{sec:retune}~节的重标定把规则一改为主推 $L=10$（备选 $20$）、规则二改为主推 $L=250$（备选 $60$），其余维度不变。“成本占毛利”一行同样来自论文口径，工程增强后的实测值见第~\ref{sec:verify}~节。

\subsection{三套规则的共同点与差异}

\textbf{唯一的形式差异是符号，其余差异全部在执行层}。这是本文（以及其依据的论文）最核心的判断。表中“信号口径”一栏看起来三种不同，但拆开看：规则一用 $+\mathrm{sign}(s)$，规则二用 $-\mathrm{sign}(s)$，规则三用 $-\mathrm{sign}(e)$，而 $e$ 本身就是两资产价格的对数差，等价于价差的累计收益。三者都是“过去价格变化的单调函数”，符号分别是 $+$、$-$、$-$。

真正的差异在四行：\emph{波动率目标化的取舍}、\emph{入场门槛的有无}、\emph{中性化的必要性}、\emph{防线的内容}。这四行恰好对应论文七组对偶中最具工程后果的四组。

\textbf{为什么规则一用波动率目标而另两套不用}。这不是风格偏好，而是信号与波动率的联合分布不同。趋势信号在波动率上升时质量下降，反转信号在高波动时反而更强（偏离更大、回归空间更大）。若对三套规则统一施加波动率目标，会系统性地削弱规则二与规则三——论文执行章称这是“在信号最强时把仓位压到最小”。

\textbf{为什么只有规则三有协整门槛}。论文对偶七指出，反转策略在结构上应带资格检验（“先验证后交易”），而趋势策略没有这道门槛（“边交易边估计”）。规则二虽然也是反转，但它走的是全市场排序，排序本身隐含了流动性筛选，没有第二道结构门槛。规则三做配对，则必须显式检验协整关系——否则“价差会回归”这一前提不成立。

\subsection{选用决策树}

\begin{enumerate}[leftmargin=2.2em]
  \item \textbf{能否做空？}
  \begin{itemize}[leftmargin=1.8em]
    \item 不能（A 股账户为主，无期货或融券）：只能选\textbf{规则二}，且必须走版本 A（纯多头 + 指数对冲）。
    \item 能（有期货账户）：继续。
  \end{itemize}
  \item \textbf{期望持仓周期？}
  \begin{itemize}[leftmargin=1.8em]
    \item 少于 5 个交易日：三套规则都不适用。论文实证显示短窗口在全部口径上都被成本摧毁（期货趋势 $L=5$ 成本占毛利 $50.1\%$；A 股反转 $L=5$ 毛利夏普 $1.08$ 被压成 $-0.35$）。
    \item 5 至 60 个交易日：继续。
  \end{itemize}
  \item \textbf{首要目标？}
  \begin{itemize}[leftmargin=1.8em]
    \item 低换手、大资金：选\textbf{规则一}的 $L=250$ 档（换手 $8.83$，成本占毛利 $10.2\%$）。
    \item 追求最高历史夏普：选\textbf{规则二}的 $L=60$ 档（净夏普 $0.99$），前提是接受 A 股做空实现与除权修复的残余偏差。
    \item 已有期货账户、希望降低商品指数方向暴露：选\textbf{规则三}，但须自行回测配对版本的收益水平。
  \end{itemize}
  \item \textbf{是否需要同时运行多套？}
  \begin{itemize}[leftmargin=1.8em]
    \item 若要同时运行规则一与规则三，必须注意论文对偶六的警告：两策略的失效在极端去杠杆事件中会同时发生（2007 年 8 月），因为冲击来自共同的资金端而非各自的信号端。此时“分散化”会失效，组合层面须预留额外冗余。
  \end{itemize}
\end{enumerate}

\emph{更正（见第~\ref{sec:retune}~节）}：第 3 步的窗口取值沿用论文口径，工程增强后的重标定结果与之不同——规则一首选 $L=10$、规则二首选 $L=250$。决策树的分支逻辑不变，仅窗口的具体取值以第~\ref{sec:retune}~节为准。

\subsection{一个必须避免的组合方式}

有一种做法看起来合理但实际有害：\textbf{把趋势与反转等权叠加，期望获得分散化}。论文执行章对此的分析是，两策略在某些状态下相关性趋近于一（共同的资金冲击），在另一些状态下相关性为负（信号符号相反）；简单叠加的实际效果是“在信号相反时互相抵消，在资金冲击时同时亏损”。

正确的融合方式有三条（均为工程建议，论文未检验）：
\begin{itemize}[leftmargin=2.2em]
  \item \textbf{按尺度分层}：短尺度配反转、中尺度配趋势，让不同频段各自发挥作用；
  \item \textbf{按状态加权}：用 regime 概率动态调整两策略的权重，而非固定比例；
  \item \textbf{按风险预算而非资金分配}：先评估两策略在目标状态下的波动与尾部，再分配风险预算。
\end{itemize}
